For an easier handling of transformers, we define a notion of normalized transformer.
This notion will specify properties of the $\OUTPUT$ matrix and of the vector consumed by this step.

\begin{definition}[Normalized transformer]
    We will say a transformer is normalized when it meets two conditions. First, the matrix of the $\OUTPUT$ has to be a projection, i.e. it's a diagonal matrix with only ones in its diagonal. Second, the last vector after the last iteration has to only be composed of $\minrep$ in every projected coordinate except, one which has to be $\maxrep$. The non projected coordinates are free to be any value.
\end{definition}

The motivation for this definition is to simplify the $\OUTPUT$ layer (first constraint) and to have more control over the probability distribution generated (second constraint). The exact values of the last vectors are such that after projecting and applying $\softmax$ in the $\OUTPUT$ layer we get a deterministic distribution with $100\%$ of the probability concentered in only one token.



\subsection*{Algorithm for normalizing} {\color{orange} not sure about title}

Let's see how to normalize any transformer.
Let's take a transformer and modify it to meet the first constraint.

\smallskip

Every linear transformation $M: \Rd \to \mathbb{R}^{|\Sigma|}$ with $d \geq |\Sigma|$ can be decomposed in two steps. \footnote{We can always add more coordinates to the embedding if $d < |\Sigma|$. Since the language is fixed, it will be a constant number of coordinates, thus not changing the class of the transformer.}

\[
    \left( \begin{matrix} x_1    \\ \vdots \\ x_d    \end{matrix} \right) 
    \xrightarrow{M'}
    \left( \begin{matrix} (Mx)_1 \\ \vdots \\ (Mx)_{|\Sigma|} \\ v \end{matrix} \right)
    \xrightarrow{projection}
    \left( \begin{matrix} (Mx)_1 \\ \vdots \\ (Mx)_{|\Sigma|} \end{matrix} \right)
\]

\noindent where $v \in \mathbb{R}^{d-|\Sigma|}$ can have any values.

Let's $\TF$ be a transformer and $M$ be the matrix of the $\OUTPUT$ layer.
Using the decomposition just presented, and the strategy defined in the previous section for applying linear transformations, we can change the matrix of the $\OUTPUT$ layer to be only a projection.
The linear transformation has to be applied to the last vector.

\medskip

Let's now make it meet the second condition.

Once $\TF$ has been modified to meet the first condition, we can use the maximum operation described in the previous section to make all the positions of the vector $0$ except the one that was the maximum before.
Now, if the maximum is positive, multiplying by $\maxrep$, then subtracting $1$ to every coordinate and multiplying by $\maxrep$ again gets us what we wanted.

\[
    \left( \begin{matrix} 0 \\ \vdots \\ 0 \\ x_{max} \\ 0 \\ \vdots \\ 0    \end{matrix} \right) 
    \xrightarrow[\maxrep]{\text{multiplying by}}
    \left( \begin{matrix} 0 \\ \vdots \\ 0 \\ \maxrep \\ 0 \\ \vdots \\ 0    \end{matrix} \right) 
    \xrightarrow{\text{substracting } 1}
    \left( \begin{matrix} -1 \\ \vdots \\ -1 \\ \maxrep -1 \\ -1 \\ \vdots \\ -1    \end{matrix} \right) 
    \xrightarrow[\maxrep]{\text{multiplying by}}
    \left( \begin{matrix} \minrep \\ \vdots \\ \minrep \\ \maxrep \\ \minrep \\ \vdots \\ \minrep    \end{matrix} \right) 
\]

It could be the case that the maximum was negative, in that case the resulting vector won't be what we are looking for.
Let's see how to make the maximum strictly positive so we can add this mechanism in the middle.

By this point our vector looks like $(0, \dots, 0, x_{max}, 0, \dots, 0)$, we are going to transform it to $(0, \dots, 0, |x_{max}|, 0, \dots, 0)$.
Remember that the coordinates we are interested in are only the first $|\Sigma|$.
Let's call $v$ to the block containing these coordinates.
Our algorithm will do the following:

\[
    \left(\begin{matrix}  v \\ \vdots \end{matrix}\right)
    \xrightarrow[]{}
    \left(\begin{matrix}  v \\ -v \\\vdots \end{matrix}\right)
    \xrightarrow[]{}
    \left(\begin{matrix}  \relu(v) \\ \relu(-v) \\\vdots \end{matrix}\right)
    \xrightarrow[]{}
    \left(\begin{matrix}  \relu(v) + \relu(-v) \\ \vdots \end{matrix}\right)
\]

\noindent where $\relu$ is applied position wise. 
This works since for all $x$ it holds that $\relu(x) + \relu(-x) = |x|$.

The first step can be done in the same way as the first half of the linear transformations.
For the second step we can take max coordinate in $v$ and $-v$.
This works because all the coordinates are $0$ except one, which if it is negative will become $0$ and if it is non-negative will remain the same.
For adding both vectors we can use one $\FF$ with the following parameters.
Notice that this works since all the coordinates of $\relu(-v)$ are non-negative.


\begin{align*}
    & W_1 = id^{d\times d}
    & W_2  = \left(\begin{matrix}
        0^{|\Sigma| \times |\Sigma|}    & id^{|\Sigma| \times |\Sigma|}   & 0^{|\Sigma| \times (d - 2*|\Sigma|)}\\
        0^{d-|\Sigma| \times |\Sigma|}  & 0^{d-|\Sigma| \times |\Sigma|}  & 0^{d-|\Sigma| \times d-2*|\Sigma|} \\
    \end{matrix}\right) &
    & b_1 = b_2 = \left(\begin{matrix}
        0^{d\times d}
    \end{matrix}\right) 
\end{align*}
